\subsection{Frobenius Reciprocity}

Let A be a semisimple ring. Let f be a homomorphism from A to a semisimple ring B. Let S be a simple A-module and T a simple B-module, and let D and E be the centralizers of S and T, respectively. By Schur's lemma (VIII, p.~47, Corollary), D and E are fields. Let H be the set of A-linear homomorphisms from S to \(f_{*}(\mathrm{T})\) . We endow H with an (E,D)-bimodule structure with external laws \((e,u)\mapsto e\circ u\) and \((d,u)\mapsto u\circ d\) for \(e\in E\) , \(u\in H\) , \(d\in D\) .

PROPOSITION 11. --- a) The multiplicity \([f_{*}(\mathrm{T}):\mathrm{S}]\) of the simple A-module S in the semisimple A-module \(f_{*}(\mathrm{T})\) is equal to the dimension of H viewed as a right vector space over D.

\begin{enumerate}
\def\labelenumi{\alph{enumi})}
\setcounter{enumi}{1}
\tightlist
\item
  The multiplicity \([f^{*}(\mathrm{S}):\mathrm{T}]\) is finite and equal to the dimension of H viewed as a left vector space over E.
\end{enumerate}

Assertion a) follows from formula (11) of VIII, p.~72.

The B-module \(f^{*}(S)\) is semisimple and finitely generated, hence has finite length. By formula (12) of loc. cit., we have

\[
[ f ^ {*} (\mathrm{S}): \mathrm{T} ] = \dim_ {\mathrm{E}} \operatorname{Hom} _ {\mathrm{B}} (f ^ {*} (\mathrm{S}), \mathrm{T}).\tag{20}
\]

Now, in II, §5, No.~1, p.~277-278, formula (2) and Remark 2, we defined an E-linear bijection from \(\mathrm{Hom}_{\mathrm{A}}(\mathrm{S},f_{*}(\mathrm{T}))\) to \(\mathrm{Hom}_{\mathrm{B}}(f^{*}(\mathrm{S}),\mathrm{T})\) . Assertion b) then follows from formula (20).

COROLLARY (Frobenius reciprocity). --- Suppose that A and B are semisimple algebras that are finite-dimensional over a commutative field K and that f is K-linear. Then the K-vector spaces S, T, D, E, and H are finite-dimensional, and we have the equalities

\[
[ f _ {*} (\mathrm{T}): \mathrm{S} ] [ \mathrm{D}: \mathrm{K} ] = [ f ^ {*} (\mathrm{S}): \mathrm{T} ] [ \mathrm{E}: \mathrm{K} ] = [ \mathrm{H}: \mathrm{K} ].\tag{21}
\]

In particular, when K is algebraically closed, we have D = E = K and

\[
[ f _ {*} (\mathrm{T}): \mathrm{S} ] = [ f ^ {*} (\mathrm{S}): \mathrm{T} ] = [ \mathrm{H}: \mathrm{K} ].\tag{22}
\]

Since A-module S is simple, it is monogenous, and D is a linear subspace of \(\operatorname{Hom}_{\mathrm{K}}(\mathrm{S},\mathrm{S})\) ; hence S and D are finite-dimensional over K. For an analogous reason, T and E are finite-dimensional over K. Finally, H is a linear subspace of \(\operatorname{Hom}_{\mathrm{K}}(\mathrm{S},\mathrm{T})\) ; it is therefore also finite-dimensional. Formula (21) then follows from Proposition 11 because the dimension of H over K is equal to \([H:D][D:K]\) and to \([H:E][E:K]\) (II, §1, No.~13, p.~222, Proposition 25). By Theorem 1 of VIII, p.~47, if K is algebraically closed, we have D = E = K. The second part of the corollary then follows from the first.

Let \(\mathrm{A}\) and \(\mathrm{B}\) be semisimple algebras that are finite-dimensional over a commutative field \(\mathrm{K}\) , and let \(f\) be a homomorphism of \(\mathrm{K}\) -algebras from \(\mathrm{A}\) to \(\mathrm{B}\) .

Let \(\mathcal{S}(\mathrm{A})\) be the set of classes of simple A-modules; for every \(\lambda \in \mathcal{S}(\mathrm{A})\) , let \(\mathrm{S}_{\lambda}\) be a module of class \(\lambda\) and \(\mathrm{D}_{\lambda}\) its centralizer. Then \(\mathrm{D}_{\lambda}\) is an algebra of finite degree over \(\mathrm{K}\) ; we denote this degree by \(d_{\lambda}\) . We define \(\mathcal{S}(\mathrm{B})\) , \(\mathrm{T}_{\mu}\) , \(\mathrm{E}_{\mu}\) , and \(e_{\mu}\) for \(\mu\) in \(\mathcal{S}(\mathrm{B})\) analogously. The Grothendieck group \(\mathrm{K}_0(\mathrm{A})\) has the family \(([\mathrm{S}_{\lambda}])_{\lambda \in \mathcal{S}(\mathrm{A})}\) as a basis, and \(\mathrm{K}_0(\mathrm{B})\) has \(([\mathrm{T}_{\mu}])_{\mu \in \mathcal{S}(\mathrm{B})}\) as a basis. Let \((a_{\mu \lambda})\) be the matrix of \(f^{*}: \mathrm{K}_{0}(\mathrm{A}) \to \mathrm{K}_{0}(\mathrm{B})\) and \((b_{\lambda \mu})\) the matrix of \(f_{*}: \mathrm{K}_{0}(\mathrm{B}) \to \mathrm{K}_{0}(\mathrm{A})\) with respect to these bases. By definition, we have

\[
a _ {\mu \lambda} = [ f ^ {*} (\mathrm{S} _ {\lambda}): \mathrm{T} _ {\mu} ], \qquad b _ {\lambda \mu} = [ f _ {*} (\mathrm{T} _ {\mu}): \mathrm{S} _ {\lambda} ]\tag{23}
\]

for \(\lambda\) in \(\mathcal{S}(\mathrm{A})\) and \(\mu\) in \(\mathcal{S}(\mathrm{B})\) . We denote the dimension of the vector space \(\mathrm{Hom}_{\mathrm{A}}(\mathrm{S}_{\lambda}, f_{*}(\mathrm{T}_{\mu}))\) over the field K by \(h_{\lambda \mu}\) . By the corollary above, we have

\[
h _ {\lambda \mu} = e _ {\mu} a _ {\mu \lambda} = d _ {\lambda} b _ {\lambda \mu}.\tag{24}
\]

When the field K is algebraically closed, we have \(d_{\lambda} = e_{\mu} = 1\) ; consequently, we have

\[
a _ {\mu \lambda} = b _ {\lambda \mu} = h _ {\lambda \mu}.\tag{25}
\]

In other words, the matrices of \(f_{*}\) and \(f^{*}\) with respect to the given bases of \(\mathrm{K}_0(\mathrm{A})\) and \(\mathrm{K}_0(\mathrm{B})\) are each other's transposes.

